\subsection{From synthetic priors to explicit autoregressive mechanisms}
\label{sec:generator-ar-motivation}

Synthetic pretraining motivates a generator that combines diversity with a
retained description of how each series was produced. KernelSynth, introduced
with Chronos, samples sums and products of Gaussian-process kernels and then
draws complete trajectories from the resulting covariance
\citep[Section~4.2]{ansari2024chronos}. Chronos-2 uses it alongside AR, ETS, and
TSI base generators before introducing multivariate dependence
\citep{ansari2025chronos2}. TimePFN's LMC-Synth mixes independent KernelSynth
latent trajectories through a shared linear mixing matrix
\citep{taga2025timepfn}. Toto~2.0 uses the separate \emph{TempoPFN} prior,
which combines several procedural and stochastic families
\citep{moroshan2025tempopfn,khwaja2026toto2}. Section~1.3 describes these
generation procedures and their computational organization.

An explicit autoregressive equation offers a complementary representation:
each dependence is attached to a source series, a lag, and an operation.
For a scalar linear example,
\begin{equation}
 x_t=b+\sum_{\ell=1}^{L}c_{\ell}x_{t-\ell}+\varepsilon_t.
 \label{eq:implemented-ar}
\end{equation}
The nonlinear AST generator replaces the weighted sum with an expression
containing selected lags and interactions. Here $L$ is the largest lag required
by the mechanism; a forecasting model may receive a longer observed context.

\paragraph{What the AR experiments establish.}
The executed KernelSynth prediction study evaluates 33 entries of the official
kernel bank and 20 sampled compositions. For each mechanism it draws 16 training
and 16 independent test trajectories of length 1024. A pooled linear AR model
with an intercept is fitted by ridge regression, with penalty $10^{-5}$ on
the lag coefficients. The orders tested are
$\{1,2,4,8,16,32,64,128,256\}$. Predictions use the observed preceding values
at every test position. The reported score is the pooled mean squared error
divided by the pooled test-target variance plus $10^{-12}$.

The companion TimePFN experiment evaluates 33 latent series from 11 LMC
mechanisms, using three latents per mechanism and orders up to 128. This is an
\emph{oracle-latent} experiment: the latent series are supplied directly to
the AR models. Table~\ref{tab:implemented-ar-evidence} summarizes the saved
results.\footnote{Executed records and all individual plots:
\href{https://github.com/MateusWiteck/tsfmxai/blob/27ec9cf/experiments/autoregressive_process_coverage/01_kernelsynth_ar_prediction_results.ipynb}{KernelSynth AR predictions}
and
\href{https://github.com/MateusWiteck/tsfmxai/blob/27ec9cf/experiments/autoregressive_process_coverage/02_timepfn_lmc_ar_prediction_results.ipynb}{TimePFN latent AR predictions}.
Each notebook also identifies its source experiment and reproduction command.}

\begin{table}[H]
\centering\small
\begin{tabular}{@{}lrrr@{}}
\toprule
Evaluated set & Cases & NMSE $\leq 0.001$ & Percentage \\
\midrule
KernelSynth bank and compositions & 53 & 45 & 84.91\% \\
TimePFN oracle latent series & 33 & 31 & 93.94\% \\
\bottomrule
\end{tabular}
\caption{Best observed one-step AR score over the tested orders.}
\label{tab:implemented-ar-evidence}
\end{table}

Order selection uses the smallest score on the same test set being summarized,
so these percentages describe exploratory best-of-grid approximation capacity.
They motivate explicit lagged mechanisms while leaving independent model
selection, recursive multi-step forecasting, distributional reproduction, and
recovery of the original GP mechanism as separate evaluation questions. Both
white-noise entries remain above the threshold, as do six KernelSynth
compositions. White-noise innovations contribute irreducible one-step
uncertainty. Constant trajectories also require care: pooled variance includes
differences between trajectories even when each trajectory changes very little.
